\scanpage{51}
% archive sheet number: 50
\mspage{\handwritten{13}}
\[
\handwritten{\mathscr H^0}:C_A^{\handwritten{\sim}}\handwritten{\longrightarrow}C_A^{\handwritten{\wedge}}\tag*{\handwritten{(10)}}
\]
le foncteur d'injection des faisceaux dans les préfaisceaux.

Le foncteur $\handwritten{\mathscr H^0}$ est un foncteur exact à ga\deleted{\unsure{n|u}}uche (2.3.14.II). Ses foncteurs dérivés droitsseront notés :
\[
\handwritten{\mathscr H^q}:C_A^{\handwritten{\sim}}\handwritten{\longrightarrow}C_A^{\handwritten{\wedge}}\tag*{\handwritten{(11)}}
\]
Pour tout objet $X$ de $C$ \deleted{\unsure{, on a|on a}} et pour tout $A$-Module $F$, on a :
\[
\handwritten{\mathscr H^q(F)(X)}\ \deleted{\unsure{\mathscr H^q|H^q}}\ 
\deleted{\unsure{H^0(C^{\wedge}/X,\mathscr H^q(F))|H^0(C/X,\mathscr H^q(F))}}=H^q(C^{\handwritten{\sim}}/X,F)\tag*{\handwritten{(12)}}
\]
En effet pour $q=0$, cet isomorphisme résulte immédiatement des définitions . On l'en déduit pour tout $q$ en remarquant que le foncteur\char34 section\handwritten{\inserted{s}} sur $X$\char34\ est exact \deleted{\unsure{à|a}} sur la catégorie des préfaiscea\deleted{\unsure{ux|us}}ux .

Soit $L:C^{\handwritten{\wedge}}\handwritten{\longrightarrow}C^{\handwritten{\wedge}}$ le foncteur défini en (2.3.11.II). Le foncteur $L$ commute aux limites projectives finies et par suite transforme \inserted{les} préfaisceaux d' $A$-modules en préfaisceaux de $A$-modules. Nous noterons :
\[
\handwritten{\check{\mathscr H}^{0}}:C_A^{\handwritten{\wedge}}\handwritten{\longrightarrow}C_A^{\handwritten{\wedge}}\tag*{\handwritten{(13)}}
\]
le foncteur induit par $L$ sur les préfaisceaux de $A$-modules.

Le foncteur $\handwritten{\check{\mathscr H}^0}$ est exact à \deleted{droite} gauche. Ses foncteurs dérivés droits seront notés $\handwritten{\check{\mathscr H}^q}$. \deleted{Les préfai} Soit $G$ un préfaisceau de $A$-modules. Les préfaisceaux $\handwritten{\check{\mathscr H}^q}(G)$ seront appelés les préfaisceaux de \uline{cohomologie de $\overset{\handwritten{\scriptscriptstyle\vee}}{\text{C}}$ech} du pré\deleted{g}faisceau $G$.

P\unsure{p|{}}our tout objet $X$ de $C$, on a un isomorphisme canonique :

\deleted{\[
\deleted{H^0(\unsure{C/|C^\wedge})\ \unsure{H^0(C^\wedge/X,(G))|H^0(C/X,(G))}\ \operatorname{Lim}}
\]
\[
H^0(C^{\wedge}/X,\quad(G))=\mathop{\mathrm{Lim}}_R H^q(C^{\wedge}/R,G)^{\unsure{C|G}}
\]
En effet pour $q=0$, cet isomorphisme résulte immédiatement des définitions}

\[
\handwritten{\check{\mathscr H}^q(G)(X)}\,\deleted{\unsure{H^0(C^\wedge/X,\check{\mathscr H}^q(G))|H^0(C/X,\mathscr H^q(G))}}=
\mathop{\underset{\handwritten{\longrightarrow}}{\mathrm{Lim}}}_R H^q(C^{\handwritten{\wedge}}/R,G)\tag*{\handwritten{(14)}}
\]
\deleted{Ainsi qu'}\inserted{\handwritten{Comme}} il resulte immédiatement des définitions (2.3.11.II).

(Dans cette formule la limite inductive est prise suivant la catégorie des cribles couvrant\deleted{s} $X$).
